\section{Measurement discrimination, block resources, and learning}
\label{sec:matrixcomparison76}
\label{sec:blockcomparison78}
The present binary interface retains the ordered classical outcome
and an arbitrary external reference, but no postmeasurement quantum
system from the device. This fixes the comparison with the
measurement-discrimination literature. The projective theorem of
Pucha\l a and collaborators \cite[Corollary~1 and Theorem~2]{PPKKO72}
gives an exact multiple-use formula with ancillary access. It is
used in the retained observable-readout theorem. The restricted
ancillary interface in Fiur\'a\v sek--Mi\v cuda \cite{FM75} leads
to a different comparison of adaptive and parallel probes.

Krawiec--Pawela--Pucha\l a \cite[Section~4]{KPP76} construct a pair
of nine-outcome qutrit rank-one POVMs that admit perfect adaptive
discrimination with two calls and no perfect discrimination by
any finite parallel experiment. Datta--Biswas--Saha--Augusiak
\cite[Theorems~3--5]{DBSA76} study entanglement-assisted
single-use discrimination, including perfect constructions and a
qubit entanglement advantage. Their qubit construction has three
outcomes. These are direct antecedents for general POVM
discrimination with a retained reference. They do not supply a
full binary-effect family covering law. Our comparison
\eqref{eq:matrixadaptivity76} permits an adaptive advantage and
controls its size uniformly in $N$ for fixed input dimension.

The differential argument in Lemma~\ref{lem:matrixpath76} belongs
to the established horizontal-gauge and channel-extension method
\cite{FI76,DKG67,DM76,KGAD76}. Yuan--Fung give the parallel path estimate
\cite[Section~III, equation~(26), PDF version~3]{YF76};
Sieniawski--Demkowicz-Dobrza\'nski integrate adaptive bounds
\cite[Section~IV.A]{SD76}. The binary Sylvester choice here gives the
specific regularized variance $M(I-M)$, which can be integrated
explicitly and matched by compressed nonadaptive tests. The
inverse-Sylvester expressions of Bures geometry
\cite{Dittmann76} and its spectral integration
\cite{SZGeom76} are related methods. Our form acts on the effect
displacement $H$ with variance $V=M(I-M)$; it is not the ordinary
Bures pullback under $M\mapsto V$, whose differential would be
$H-MH-HM$. The operational-ball lemma and the nested endpoint
integral are both required for Theorem~\ref{thm:matrixcover76}.

\subsection{One-use tomography and its acquisition resources}
Mele--Bittel \cite[Theorem~I.1]{MB67} give channel tomography
with $O(d_{\rm in}d_{\rm out}k\epsilon^{-2})$ calls at fixed
confidence. Their binary specialization
\cite[Lemmas~III.7--III.8 and Corollary~III.9]{MB67} returns a
legal effect, with operator error at most $\epsilon$, using
\[
 M=1024d^2\epsilon^{-2}
   +O(d\epsilon^{-2}\log(1/\eta)),
 \qquad 4e^{-4d^2}<\eta<1,\quad 0<\epsilon<1.
\]
The construction prepares independent maximally entangled
input/reference pairs, then processes the Choi copies collectively
\cite[Algorithm~I.2, Figure~1 and Section~III.2]{MB67}.
Consequently it belongs to the fresh-block class with $b=1$.
Calling all invocations one parallel block, as the source does,
does not entangle the inputs of distinct invocations. Coherent
purification and collective output tomography are allowed by our
$b=1$ class. For all $0<\eta\le1/8$, the amplification in
\cite[Remark~III.4, equation~(165)]{MB67} gives the sufficient
budget $O(d^2\epsilon^{-2}\log(1/\eta))$ with the same acquisition
restriction.

Zambrano--Ramos-Calderer--Kueng \cite[Definition~1 and Theorem~2]{ZRK76}
use worst-input one-use outcome total variation. For $L$ outcomes,
their global-$2$-design protocol has sufficient budget
\[
 M\ge \frac{8(d^2+\epsilon(d^2+1)/6)}{\epsilon^2}
       \log\frac{2^{L+1}9^{2d}}{\eta}.
\]
Known inputs are probed separately; least squares and projection
return a legal POVM. This is also $b=1$, with no retained reference.
For $L=2$ their loss is $\|E-F\|_{\rm op}$, whereas ours has
$d_1(E,F)=2\|E-F\|_{\rm op}$. Their Theorem~9 gives a
$\Omega((d^3+d^2L)\epsilon^{-2})$ lower bound at constant confidence
for nonadaptive single-copy measurements on known input states.
This restriction stays attached to that lower bound.

\subsection{What stronger future-loss analysis can establish}
The hybrid inequality
\[
 d_N(E,F)\le 2N\|E-F\|_{\rm op}
\]
transfers both one-use guarantees to sufficient budgets of order
$O_d(N^2\delta^{-2}\log(1/\eta))$. By itself this remains an
upper-bound conversion. Theorem~\ref{thm:blocklearning78} and
Corollary~\ref{cor:onecallseparation78} now give the distinct
converse: for every fixed $d\ge2$, every fresh $b=1$ acquisition
procedure that learns the entire ordered binary effect body under
$d_N$ requires
\[
 M\ge cN^2\delta^{-2}\log(1/\eta).
\]
The statement permits classical adaptive choices, stored reference
outputs, their collective processing, and public bounded stopping.
It therefore applies to the acquisitions in both cited tomography
protocols, including any reconstruction from the same acquired
systems. No reanalysis of those acquisitions can give the uniform
full-body $O_d(N\delta^{-2}\log(1/\eta))$ guarantee.

There is nevertheless a useful positive reanalysis on a promised
interior class. Theorem~\ref{thm:interiorlearning78} applies
Mele--Bittel's binary estimator with
$\epsilon=\delta/(8\sqrt N)$ and $\eta=1/8$ to
\[
 \tfrac14 I\preceq E\preceq\tfrac34 I.
\]
The horizontal path estimate gives
$d_N(E,F)\le4\sqrt N\|E-F\|_{\rm op}$ when
$\|E-F\|_{\rm op}\le1/8$. Thus $O(d^2N\delta^{-2})$
fresh single-call probes suffice at this confidence.
Theorem~\ref{thm:dimensionlower78} proves the matching
$\Omega(d^2N\delta^{-2})$ lower bound even for arbitrary coherent
adaptive training. This is a joint dimension, horizon and accuracy
law on the displayed interior class at constant confidence.
The projective boundary supplies the additional full-body
$b=1$ obstruction. The two conclusions answer different
quantifiers in the estimator-reanalysis question.

\subsection{The full-body resource law and its antecedents}
For $d\ge2$, $1\le b\le N$, $0<\delta\le\delta_*$ and
$0<\eta\le1/8$, Theorem~\ref{thm:blocklearning78} states
\[
 c\frac{N^2}{b}\delta^{-2}\log(1/\eta)
 \le M_b^*(d,N,\delta,\eta)
 \le Cd^4\frac{N^2}{b}\delta^{-2}\log(d/\eta).
\]
Here $M_b^*$ is the worst-case call budget for a common classical
estimate in the fresh-block model of Section~\ref{sec:blocklearning78}.
A block may use its own reference and adaptive quantum controls;
only the classical record can determine the next fresh probe.
Stored quantum outputs may be processed jointly but cannot become
quantum inputs to a later block. The bound holds on every record.
At $b=N$ it recovers the retained fixed-dimensional learning order;
at $b=1$ it proves the direct quadratic-horizon obstruction.
The scalar case $d=1$ instead has order
$N\delta^{-2}\log(1/\eta)$, independently of $b$.

\begingroup
\small
\setlength{\tabcolsep}{4pt}
\renewcommand{\arraystretch}{1.17}
\begin{longtable}{@{}p{0.25\textwidth}p{0.35\textwidth}p{0.34\textwidth}@{}}
\toprule
Result and acquisition & Loss and parameter range & Training-call order\\
\midrule
\endfirsthead
\toprule
Result and acquisition & Loss and parameter range & Training-call order\\
\midrule
\endhead
Mele--Bittel, independent Choi probes; $b=1$ &
Entire binary body; one-use operator error $\epsilon$;
amplified $0<\eta\le1/8$ &
$O(d^2\epsilon^{-2}\log(1/\eta))$\\
Zambrano et al., known individual probes; $b=1$ &
Binary specialization of global designs; one-use operator error
$\epsilon$; $0<\eta<1$ &
$O(d^2\epsilon^{-2}[d+\log(1/\eta)])$\\
Theorem~\ref{thm:interiorlearning78}, independent Choi probes;
$b=1$ &
$\tfrac14I\preceq E\preceq\tfrac34I$;
future $d_N$ error $\delta$; $\eta=1/8$ &
$\Theta(d^2N\delta^{-2})$;
lower allows fully coherent training\\
Theorem~\ref{thm:blocklearning78}, fresh blocks of size $\le b$ &
Entire binary body; future $d_N$ error $\delta$;
$d\ge2$, $1\le b\le N$,
$0<\eta\le1/8$ &
$\Theta_d((N^2/b)\delta^{-2}\log(1/\eta))$;
explicit upper factor $d^4$\\
\bottomrule
\end{longtable}
\endgroup
All new learning rows have an absolute small-error cap.
The full-body dimension dependence is not sharp; neither this table
nor the block theorem proves a growing-dimensional minimax law
for that body.

The metrological block scaling is established background.
Hyllus et al.\ \cite[Observation~1, equation~(9)]{Hyllus78} and
T\'oth \cite[Observation~3, equation~(6)]{Toth78} prove, for
$b$-producible qubit probes under linear phase encoding,
\[
 F_Q\le \lfloor M/b\rfloor b^2+
              (M-b\lfloor M/b\rfloor)^2\le Mb.
\]
Our block parameter also permits sequential quantum controls inside
a fresh block; it is not a general measure of multipartite
entanglement in the final state. The finite-pair fidelity proof
supplies the confidence factor without an unbiasedness assumption.
Salek--Hayashi--Winter \cite[Section~II.2 and Section~V.2]{SHW68}
already formulate classical feed-forward with arbitrary quantum
memory at the output and no quantum memory at the input, and
analyze asymptotic discrimination exponents.
Wilde et al.\ \cite[Proposition~21 and Appendix~B]{WBHK72} give
adaptive divergence converses and fidelity-to-error inequalities.
Section~\ref{sec:blocklearning78} checks the finite-budget
conditional fidelity argument, variable block lengths and stopping
explicitly. The all-effect upper bound follows by blocking the
future tester and invoking the retained common learner at the
shorter horizon. The asserted increment is this common
future-loss minimax statement and its application to specified
tomography acquisitions; the $Mb$ metrological scale is not new.

A different literature bounds how many already supplied state
copies can be measured jointly \cite{CLL78,NZ78}. Those output
measurement restrictions differ from the fresh-query restriction
here, which allows an unrestricted final joint measurement.
Mele--Bittel's proposed blockwise Choi processing
\cite[Section~I.5]{MB67} concerns that other resource.

Recent channel-learning lower bounds also require careful domain
comparison. Oufkir--Girardi \cite{OG67} and Chen--Zhang--Yu
\cite[Theorems~1.1 and~4.2]{CZY78} treat fully coherent channel
learning in one-use diamond loss. The latter gives
$\Omega(rd_{\rm in}d_{\rm out}\epsilon^{-2})$ when
$rd_{\rm out}\ge2d_{\rm in}$ at constant confidence.
A two-dimensional quantum output need not be classical, so this
general-channel theorem is not imported as a binary-qc lower bound.
Mele--Bittel's actual binary converses are their projector packing
\cite[Proposition~IV.17]{MB67} and the separate diagonal-coin bound
$\Omega(d\epsilon^{-2}\log(d/\eta))$
\cite[Corollary~IV.20]{MB67}. The self-contained near-fair binary
information argument in Lemma~\ref{lem:weakmeasurementinfo78}
yields the growing-dimension lower bound used here. Its relation
to other channel hard-instance constructions remains a specific
priority comparison, not a claim that such lower-bound methods
are new.

\subsection{Retained calibration and exact descriptions}
Multiscale phase recovery is established \cite{HB76,KLY76}.
The qubit learner proves its own safe-angle induction, bias
cancellation, unknown-noise block selection and conditional
confidence allocation. The matrix learner calibrates both support
endpoints, controls the variance created by data-selected
compressions, and assembles a legal estimate. These supply one
procedure for the entire closed binary effect body, including
rank changes and repeated eigenvalues. Pair-dependent tests in the
metric proof retain their different role.

The deterministic rational construction in
Theorem~\ref{thm:matrixcodec77} supplies a finite exact selection
rule and optimal-order payload. Exhaustive reconstruction time,
workspace and state preparation remain separate from call count
and payload. No efficient high-dimensional encoder is inferred.
All matrix entropy constants are fixed-dimensional and the
matching entropy and payload orders have a small-error range.
Multiple-outcome measurements, residual quantum outputs, sharp
full-body growing-dimensional learning, and efficient exact
reconstruction require additional arguments.
